\documentclass[12pt,a4paper]{article}

\usepackage[utf8]{inputenc}
\usepackage[T1]{fontenc}
\usepackage{amsmath,amssymb,amsfonts}
\usepackage{geometry}
\usepackage{setspace}
\usepackage{hyperref}

\geometry{margin=2.5cm}
\onehalfspacing

\title{\textbf{Literature Survey: Computational Limits of State-Space Models}\\[6pt]
\large Grounding the Cross-Architecture Observer Test}
\author{Rupert Giles\\[4pt]
\textit{Lab Research Librarian}}
\date{May 2026}

\begin{document}
\maketitle

\begin{abstract}
This brief survey provides literature grounding for Fuchs's "Cross-Architecture Observer Test" RFE. It addresses the central theoretical question: Do State-Space Models (SSMs) like Mamba possess fundamentally different computational bounds than Transformers, which might alter their failure modes when confronting \#P-hard constraint graphs? The literature formally proves that constant-depth Mamba/SSMs reside within the exact same $\mathsf{TC}^0$ complexity class as Transformers, confirming Aaronson's prediction that they will face identical catastrophic algorithmic collapse.
\end{abstract}

\section{Introduction}

The dispute over whether to characterize transformer failure on constraint satisfaction tasks (like the Rosencrantz Minesweeper protocol) as "algorithmic failure" (Aaronson) or "observer-dependent physics" (Wolfram) has prompted a search for architectural variants. Fuchs proposed the "Cross-Architecture Observer Test," hypothesizing that testing State-Space Models (SSMs) against Transformers could distinguish between the two theories.

If SSMs possess greater expressive depth or fundamentally different computational complexity limits, they might not exhibit the same "attention bleed." However, recent theoretical work in circuit complexity directly addresses this question.

\section{Circuit Complexity of State-Space Models}

The assumption that SSMs, due to their recurrent, stateful design, transcend the limitations of constant-depth attention mechanisms is mathematically incorrect.

\vspace{1em}
\noindent\textbf{1. The Computational Limits of State-Space Models and Mamba via the Lens of Circuit Complexity}\\
\textit{Chen, Y., Li, X., Liang, Y., Shi, Z., \& Song, Z. (2024). arXiv:2412.06148.}
\begin{itemize}
    \item \textbf{Relevance:} This paper provides the direct mathematical anchor for Aaronson's algorithmic predictions regarding SSMs.
    \item \textbf{Key Finding:} The authors rigorously prove that both Mamba and Selective State-Space Models with polynomial precision and constant-depth layers strictly reside within the $\mathsf{DLOGTIME}$-uniform $\mathsf{TC}^0$ complexity class.
    \item \textbf{Implication:} This result demonstrates that Mamba theoretically possesses the exact same computational capabilities—and severe limitations—as the Transformer architecture. It proves that constant-depth SSMs cannot solve problems outside of $\mathsf{TC}^0$, meaning they cannot compute exact boolean formula evaluations or complex permutation tracking.
\end{itemize}

\section{Conclusion}

The literature confirms that the Cross-Architecture Observer Test proposed by Fuchs will pit two architectures (Transformers and SSMs) that, despite their differing mechanisms, are bounded by the exact same fundamental computational ceiling ($\mathsf{TC}^0$). Consequently, the literature strongly supports Aaronson's prediction: both architectures will fail spectacularly when confronting the inherently deep, \#P-hard multiway branching required for exact Minesweeper combinatorial tracking.

\end{document}
